\documentclass[11pt]{article}

\usepackage[a4paper,margin=0.88in]{geometry}
\usepackage{amsmath,amssymb,amsthm,mathtools}
\usepackage[hidelinks]{hyperref}
\usepackage{microtype}
\hypersetup{pdftitle={Infinitely Many Integer Solutions of y2+x2y+xz2-y+1=0}}

\newtheorem{theorem}{Theorem}
\theoremstyle{remark}
\newtheorem{remark}{Remark}

\title{Infinitely Many Integer Solutions of\\[-1mm]
\texorpdfstring{$y^2+x^2y+xz^2-y+1=0$}{y2+x2y+xz2-y+1=0}}
\author{}
\date{}

\begin{document}
\maketitle
\vspace{-2.2em}

\begin{abstract}
We give a completely explicit infinite family of integer solutions of
\[
 y^2+x^2y+xz^2-y+1=0.
\]
The construction fixes $x=-201$, converts the equation into an indefinite Pell-type norm equation, and iterates one integral norm-preserving transformation.  The proof itself is a direct calculation with a binary quadratic form; a final section explains how the seed and the transformation were found.
\end{abstract}

\section{The explicit family}

Set
\[
 a=515095,\qquad b=36332,\qquad D=201,
\]
and define integer sequences $(W_n)_{n\geq 0}$ and $(Z_n)_{n\geq 0}$ by
\begin{equation}\label{eq:recurrence}
 \begin{aligned}
 W_0&=20640, & Z_0&=299,\\
 W_{n+1}&=aW_n+DbZ_n,
 &Z_{n+1}&=bW_n+aZ_n.
 \end{aligned}
\end{equation}
Thus $Db=7302732$.

\begin{theorem}\label{thm:main}
For every $n\geq 0$ and every $\sigma\in\{1,-1\}$,
\begin{equation}\label{eq:family}
 (x,y,z)=\bigl(-201,\,W_n-20200,\,\sigma Z_n\bigr)
\end{equation}
is an integer solution of
\begin{equation}\label{eq:original}
 y^2+x^2y+xz^2-y+1=0.
\end{equation}
The triples in \eqref{eq:family}, indexed by $(n,\sigma)$, are pairwise distinct.  In particular, \eqref{eq:original} has infinitely many integer solutions.
\end{theorem}

\begin{proof}
The two numerical identities on which the construction rests are
\begin{equation}\label{eq:numerical-identities}
 a^2-Db^2=515095^2-201\cdot36332^2=1
\end{equation}
and
\begin{equation}\label{eq:initial-norm}
 W_0^2-DZ_0^2
 =20640^2-201\cdot299^2
 =20200^2-1.
\end{equation}
For arbitrary integers $W,Z$, put
\[
 W'=aW+DbZ,
 \qquad
 Z'=bW+aZ.
\]
A direct expansion, with the mixed terms cancelling, gives
\begin{align*}
 (W')^2-D(Z')^2
 &= (aW+DbZ)^2-D(bW+aZ)^2\\
 &= (a^2-Db^2)(W^2-DZ^2).
\end{align*}
By \eqref{eq:numerical-identities}, the transformation therefore preserves the quadratic form $W^2-DZ^2$.  Applying this observation inductively to \eqref{eq:recurrence}, and using \eqref{eq:initial-norm}, yields
\begin{equation}\label{eq:invariant}
 W_n^2-201Z_n^2=20200^2-1
 \qquad(n\geq 0).
\end{equation}

Now take $x=-201$, $y=W_n-20200$, and $z=\sigma Z_n$.  Since
\[
 201^2-1=40400=2\cdot20200,
\]
the left-hand side of \eqref{eq:original} becomes
\begin{align*}
 &(W_n-20200)^2+(201^2-1)(W_n-20200)-201Z_n^2+1\\
 &\qquad =W_n^2-201Z_n^2-20200^2+1=0
\end{align*}
by \eqref{eq:invariant}.  Hence every triple in \eqref{eq:family} is an integer solution.

It remains only to prove that infinitely many of these solutions are distinct.  The initial values $W_0,Z_0$ are positive, and every coefficient in \eqref{eq:recurrence} is positive; hence $W_n,Z_n>0$ for all $n$.  Moreover,
\[
 W_{n+1}=aW_n+DbZ_n>W_n,
\]
because $a>1$ and $DbZ_n>0$.  Thus $(W_n)$ is strictly increasing.  Solutions arising from different indices have different $y$-coordinates, while the two signs at a fixed index are distinct because $Z_n>0$.  Hence all triples indexed by $(n,\sigma)$ are pairwise distinct.  This proves the theorem.
\end{proof}

The first two levels of the family are
\[
 (-201,440,\pm299)
\]
and
\[
 (-201,12815057468,\pm903905885).
\]
Equivalently, the recurrence is encoded by
\begin{equation}\label{eq:closed-form}
 W_n+Z_n\sqrt{201}
 =(20640+299\sqrt{201})
  (515095+36332\sqrt{201})^n.
\end{equation}

\section{How the construction was found}

The point of the construction is not the large numerical coefficients; it is the sign of $x$ and the norm structure that this sign exposes.

\subsection*{1. Forcing an indefinite Pell-type equation}
Write $x=-A$, where $A$ is a positive odd integer, and set
\[
 H=\frac{A^2-1}{2}.
\]
Then \eqref{eq:original} is
\[
 y^2+(A^2-1)y-Az^2+1=0,
\]
so completing the square gives the equivalent equation
\begin{equation}\label{eq:general-pell}
 (y+H)^2-Az^2=H^2-1.
\end{equation}
This is indefinite.  Once one solution is available, multiplication by a nontrivial solution of $u^2-Av^2=1$ can preserve the right-hand side and generate further solutions.  This explains why a negative value of $x$ is natural.  For comparison, fixing $x=A>0$ would instead produce a positive-definite equation of the form
\[
 (y+H)^2+Az^2=H^2-1,
\]
which has only finitely many integral possibilities for that fixed $A$.

\subsection*{2. Engineering one seed}
For $x=-A$, the original expression can be grouped as
\begin{equation}\label{eq:seed-grouping}
 y^2-y+1+A(Ay-z^2).
\end{equation}
Consequently, it is enough to find integers $y,r,z$ satisfying
\begin{equation}\label{eq:seed-system}
 y^2-y+1=Ar^2,
 \qquad
 z^2=Ay+r^2.
\end{equation}
Indeed, substitution into \eqref{eq:seed-grouping} gives
\[
 Ar^2+A\bigl(Ay-(Ay+r^2)\bigr)=0.
\]
A convenient successful choice is
\[
 A=201,
 \qquad y=440,
 \qquad r=31,
 \qquad z=299,
\]
for which
\begin{equation}\label{eq:seed-verification}
 440^2-440+1=201\cdot31^2,
 \qquad
 299^2=201\cdot440+31^2.
\end{equation}
Thus $(-201,440,\pm299)$ is the required seed.  In the completed-square coordinate
\[
 W=y+\frac{201^2-1}{2}=y+20200,
\]
this seed is $(W,Z)=(20640,299)$, which explains the initial values in \eqref{eq:recurrence}.

\subsection*{3. Extracting the norm-one multiplier from the same seed}
The first equation of \eqref{eq:seed-system} has a useful hidden form, since
\[
 4(y^2-y+1)=(2y-1)^2+3.
\]
For $U,V\in\mathbb Q$, write
\[
 N(U+V\sqrt{201})=U^2-201V^2.
\]
A direct expansion shows that $N(\alpha\beta)=N(\alpha)N(\beta)$.  For the values in \eqref{eq:seed-verification}, the preceding identity becomes
\begin{equation}\label{eq:minus-three}
 879^2-201\cdot62^2=-3.
\end{equation}
Hence the quadratic expression
\[
 \eta=879+62\sqrt{201}
\]
has norm $-3$.  Squaring and dividing by $3$ gives the integral element
\begin{align*}
 \varepsilon=\frac{\eta^2}{3}
 &=\frac{879^2+201\cdot62^2}{3}
   +\frac{2\cdot879\cdot62}{3}\sqrt{201}\\
 &=515095+36332\sqrt{201}.
\end{align*}
Its norm is $(-3)^2/3^2=1$, which is exactly identity \eqref{eq:numerical-identities}.  Multiplying $W+Z\sqrt{201}$ by $\varepsilon$ gives
\[
 (aW+201bZ)+(bW+aZ)\sqrt{201},
\]
namely the recurrence \eqref{eq:recurrence}.  Thus the same arithmetic data that produce the seed also produce the norm-preserving operation that iterates it indefinitely.

\begin{remark}
Theorem~\ref{thm:main} proves infinitude, which is one of the two alternatives in the problem.  It does not assert that the displayed family exhausts all integer solutions of \eqref{eq:original}.
\end{remark}

\end{document}
